\section{相关工作}
\label{rw}

\subsection{CUA 任务基准}

CUA 基准已经从简单的控件级交互逐步发展到日益真实且多样的任务设置。MiniWoB++~\cite{miniwob} 等早期工作使用合成环境探索基于强化学习的网页控制；后续研究借助离线演示数据集扩展到真实网页~\cite{mind2web,webarena}，继而加入多模态感知~\cite{visualwebarena} 和桌面级通用性~\cite{osworld}。较新的基准进一步走向在线评测和真实世界的多样性：Online-Mind2Web~\cite{onlinemind2web} 表明，一旦离开受控的离线环境，已有报告中的许多进展便会消失；Mind2Web~2~\cite{mind2web2} 与 CocoaBench~\cite{cocoabench} 则把覆盖范围进一步扩展到智能体式搜索和异构真实应用。面向企业的 WorkArena 与 WorkArena++~\cite{workarena,workarenaplus} 把评测带入专业 SaaS 环境，但仍围绕单一企业平台构建，跨应用协调有限，任务跨度也较短。\datasetname 针对这一系列持续存在的缺口，覆盖六个专业领域的 23 个开源 SaaS 系统，要求真正的跨应用协调，并通过自动化验证评估平均超过 100 步的长程任务。

\subsection{CUA 环境与任务构建}

构建高质量 CUA 基准环境，需要在真实性、可控性、可验证性和可扩展性之间取得平衡。WebArena~\cite{webarena} 等人工构建方法通过具备完整执行支持的专门环境实现高保真，但每增加一个应用都需要大量人工投入，难以扩展；WebArena-Infinity~\cite{wainfinity} 等生成式方法利用 LLM 合成新环境和任务，提高了可扩展性，却可能牺牲真实已部署软件的真实性；人工标注数据集~\cite{mind2web} 提供广泛的任务多样性，但缺少在线执行和自动验证。\datasetname 采用不同路径：评测基于真实、开源且可部署的 SaaS，并通过 Builder--Challenger--Refiner 流水线生成任务；LLM 产生的候选任务由领域专家反复审查其可执行性、可验证性和专业真实性。由此，真实软件带来保真度，LLM 辅助合成带来规模，专家监督则保障质量。
